\chapter{Data and Methodology}
\label{ch:method}

\section{Data Construction}

The individual-level corpus is built from the ANES pilot waves of 2016, 2018,
2019, 2020, 2022, and 2024 (132,948 rows, 3,165 respondents, 29 candidates). The
population-level component uses the GSS waves of 2004--2022, with 2024 as the
target year. The raw ANES table contains two defects. Trailing-period duplicates
(\texttt{obama.} and \texttt{obama}, \texttt{hrc.} and \texttt{hrc}) produce
8,165 rows; social-group thermometers (\texttt{black}, \texttt{white},
\texttt{dem}, \texttt{rep}) stored in the candidate column produce 23,060 rows.
Names are normalized, candidate and group targets are separated, missing or
out-of-range scores are removed (3,588), and duplicates are dropped (8,141). The
clean candidate table contains 98,159 rows.

Each observation is represented as a triple $(v_i,c_j,y_{ij}\mid t_i)$, where
$v_i$ is the voter profile, $c_j$ the target, $y_{ij}\in[0,100]$ the thermometer,
and $t_i$ the year. Candidate identity is decoupled from attributes: each
candidate is described by a metadata vector $\mathcal{M}[c]$ containing party,
ideology, age, position, and incumbent status. Three features encode social
cognition. Party match is $+1$ for a co-partisan, $-1$ for an independent, and
$0$ otherwise. Ideological distance is the seven-point distance
$\lvert\text{lcself}_v-\text{ideology}_c\rvert\in[0,6]$. A similarity index
combines party, ideology, race, gender, and age with weights
$0.25,0.25,0.20,0.15,0.15$. Cross-pressure indicators flag same-party but
ideologically distant pairs (CP1), different-party but ideologically close pairs
(CP2), and demographically similar but different-party pairs (CP3).

\section{Evaluation Protocols}

ANES is a panel: 3,080 of 3,165 respondents appear in more than one wave, and
every respondent observed in 2022 or 2024 also appears in 2016--2020. There are
zero test-only respondents. A temporal split can therefore measure transfer to a
later wave, but not generalization to unseen individuals. The temporal protocol
(T) trains on 2016--2020 (6,516 rows), validates on 2022 (2,501), and tests on
2024 (3,421). The respondent protocol (R) splits respondents 80/20 by a hash of
the respondent identifier: 5,094 training rows and 6,014 test rows. A strict cell
(T$\times$R) contains the 226 rows from 2024 whose respondents are held out. R is
the main in-distribution test; T measures the temporal boundary.

\section{Fine-Tuning and Readouts}

The model is Qwen3-8B. LoRA is applied to all attention and MLP projections with
rank $r=8$, scaling $\alpha=32$, and dropout $0.1$. Training uses three epochs, a
learning rate of $10^{-4}$, an effective batch size of 16, and a maximum sequence
length of 512. The loss is the causal language-modeling loss on the answer tokens
only; prompt tokens are masked. This yields 21.8M trainable parameters (0.27\% of
the model), and three seeds (42, 43, 44) are used.

Two readouts are used. \emph{Generation} decodes greedily and parses the first
integer in $[0,100]$; it preserves individual variation and is the primary readout
for choice. \emph{Anchor likelihood} scores eleven anchors
$a\in\{0,10,\dots,100\}$ by their sequence log-probability, normalizes with a
softmax, and reports the expectation and the argmax. For choice, the candidate
with the highest predicted score is selected within each respondent--year group;
accuracy is reported against the rating-induced choice and against chance.

\section{Population-Level Pipeline}

For full attribute simulation (FAS), every univariate and pairwise cell is
extrapolated with a four-model ensemble,
\begin{equation}
\hat p(t^*)=\frac{1.0\hat p_{\text{poly}}+1.0\hat p_{\text{loess}}
+0.6\hat p_{\text{exp}}+0.8\hat p_{\text{lin}}}{3.4}.
\end{equation}
The predicted marginals and pairwise tables are combined with iterative
proportional fitting: starting from the independence approximation,
multiplicative updates match each pairwise table and marginal, run for 800
iterations, and are normalized with Laplace smoothing. The resulting joint
distribution is used to generate five task types: batch generation (50\%),
prior-constrained sampling (25\%), scenario variation (15\%), subgroup analysis
(7\%), and iterative correction (3\%).

\section{Metrics}

For PAS, MAE, RMSE, Pearson correlation, a relaxed accuracy (same 20-point bin or
absolute error $\le10$), rating-induced choice accuracy, and expected calibration
error (ECE) are reported. For FAS, JS and symmetric KL divergence are computed on
each marginal and each pairwise joint and averaged over attributes or pairs.
